\section{Crash Triage}

A representative crash was selected for triage and
replayed with ASan to confirm and analyze
the potential bug.


\subsection{Step 1: Reproduce}

The crash input was copied to a dedicated \texttt{triage/} directory and run
through \texttt{bin/png\_fuzz} with ASan symbolization enabled:

{\scriptsize
\begin{lstlisting}[style=terminal]
$ cp findings/default/crashes/id:000076,...\
      triage/crash_example.png
$ ASAN_OPTIONS=abort_on_error=0:\
    detect_leaks=0:exitcode=0:symbolize=1 \
  ./bin/png_fuzz triage/crash_example.png

ERRORERROR: AddressSanitizer: heap-buffer-overflow
  on address 0x506000000120
READ of size 11 at 0x506000000120 thread T0
  #0 strlen  (png_fuzz+0x4a3e6)
  #1 png_push_read_iTXt  pngpread.c:1606
  #2 png_process_data    pngpread.c:41
  #3 run_progressive_read harness.c:57
  #4 main                harness.c:247

0x506000000120 is located 0 bytes after
  64-byte region [0x...e0, 0x...120)
allocated in png_malloc <- png_process_data

SUMMARY: heap-buffer-overflow in strlen
\end{lstlisting}
}

ASan reports an out-of-bounds heap read originating from a
call to \texttt{strlen()} operating on attacker-controlled chunk data. The call chain localizes the fault to libpng's
iTXt parsing logic inside
\texttt{png\_push\_read\_iTXt()}, suggesting that malformed
text metadata is processed without sufficient boundary
validation.

\subsection{Step 2: Minimize with \texttt{afl-tmin}}

To isolate the minimal conditions required to reproduce the
fault, the original 418-byte crashing input was minimized
using \texttt{afl-tmin}:

{\scriptsize
\begin{lstlisting}[style=terminal]
$ AFL_SKIP_CPUFREQ=1 afl-tmin \
    -i triage/crash_example.png \
    -o triage/crash_min.png \
    -- ./bin/png_fuzz @@
    
[+] Read 418 bytes from 'crash_example.png'.
[+] Program exits with a signal,
    minimizing in crash mode.
[+] Block removal complete, 124 bytes deleted.
[+] Symbol minimization: 14 symbols replaced.
[+] Block removal complete, 15 bytes deleted.
[+] Block removal complete, 0 bytes deleted.

 File size reduced by : 33.25% (to 279 bytes)
 Characters simplified: 125.45%
 Number of execs done : 748
\end{lstlisting}
}

To validate reproducibility, the minimized input was executed
under the ASan-instrumented binary:

{\scriptsize
\begin{lstlisting}[style=terminal]
$ ASAN_OPTIONS=abort_on_error=0:\
    detect_leaks=0:exitcode=0:symbolize=1 \
  ./bin/png_fuzz triage/crash_min.png

ERRORERROR: AddressSanitizer: heap-buffer-overflow
READ of size 1 at 0x506000000052 thread T0
  #0 png_push_read_iTXt  pngpread.c:1581
  #1 png_process_data    pngpread.c:41
  #2 run_progressive_read harness.c:57

0x506000000052 is located 1 byte after
  49-byte region [0x...020, 0x...051)

SUMMARY: heap-buffer-overflow
  pngpread.c:1581 in png_push_read_iTXt
\end{lstlisting}
}
The minimized input continues to trigger a heap out-of-bounds
read within the iTXt parsing logic, confirming the successful reduction of the input.

\vspace{-0.2cm}
\subsection{Step 3: Diagnose and Bug Classification}

The iTXt parsing logic in \texttt{png\_push\_read\_iTXt()} processes
multiple variable-length fields stored sequentially and separated by
NUL (\texttt{0x00}) delimiters. These fields are traversed using
pointer iteration without enforcing strict bounds based on the
declared chunk length.

{\scriptsize
\begin{lstlisting}[style=terminal]
/* pngpread.c:1572 */
for (lang = key; *lang; lang++)
    /* Empty loop */ ;

/* pngpread.c:1581 */
for (lang_key = lang; *lang_key; lang_key++)
    /* Empty loop */ ;

/* pngpread.c:1606 */
text_ptr->itxt_length = png_strlen(text);
\end{lstlisting}
}

If a malformed iTXt chunk omits a terminator within the allocated
region, traversal continues beyond the heap buffer, resulting in a heap-based out-of-bounds read. The issue is similar in class to previously reported libpng
text-chunk parsing vulnerabilities (e.g., CVE-2015-8540), but stems
from unchecked NUL-terminated traversal rather than integer-related
bounds checking. No exact CVE match was identified for libpng~1.4.8.
